\section{Conclusion}
\label{sec:conclusion}

We asked whether small models that run on a laptop can improve their own harness. With Continual
Harness as-is, on our tasks, they did not: self-editing lowered success for both a weak model (32\%
to 15\%) and a strong model (92\% to 77\%) compared with a fixed harness, in 5 of 6 seeds. By our
operational definition, both models are below the capability floor for this loop and these tasks.
Our logs show small models writing wrong lessons into their own prompts and keeping edits they
cannot judge, while every edit makes the prompt longer; these observations are consistent with the
drop, but we did not test them as causes.

No single added technique removed the loss. All five together, where a larger local model writes
the edits and outside checks decide which ones stay, brought the weak model back to the
fixed-harness level and above Continual Harness as-is in every seed, although that difference is
not significant after correction. On held-out tasks the same system gave the highest success in our
study (60\% against 27\%), which is suggestive but not established with 15 runs. For the strong
model, the gate rejected almost every edit, which reduced the harm but gave no gain. This combined
system is best described as safeguarded adaptation with a stronger refiner, not self-improvement by
the small model.

Future work should test longer tasks with more experience to learn from, more held-out tasks and
seeds, gates with more reruns or a stricter rule for ties, a leave-one-out study of the techniques,
a fixed harness padded to the same length as an edited one, and more models between 1B and 5B
parameters. Our code, tasks and per-run outcomes are open, so these runs can be repeated on any
laptop.
